% =====================================================================
%  第 I 部：少数のターゲットラベルが最初に直すもの
%  文の元：paper/sections/4-part1-what-few-labels-fix.tex．表は論文の表 2
%  図の値：gain_mean.csv，gain_sign.csv，cost_curves.csv，loss_split.csv，centering_t.csv，thrshare.csv，dirrecov.csv，
%          lambda_<種類>.csv，value_curve.csv・value_points.csv（模式図）．出所は各 CSV の先頭
% =====================================================================
\setband{第 I 部　少ないラベルが最初に直すもの}

\begin{frame}{ソースを混ぜると，$K$ が小さいときは得，$K \ge 25$ ではどの種類でも損になる}
  \rl{結果：利得の総当たり実験}
  \vspace{-0.2em}
  \begin{columns}[T,onlytextwidth]
    \begin{column}{0.53\textwidth}
      \centering
      \begin{tikzpicture}
        \begin{axis}[width=7.4cm, height=5.5cm, xmin=-0.3, xmax=5.3, ymin=-6, ymax=15,
            xtick={0,...,5}, xticklabels={1,5,10,25,50,全件},
            xlabel={感情あたりのターゲットのラベル $K$}, ylabel={利得（UAR の\pt）},
            ymajorgrids, grid style={mMuted!25},
            legend style={font=\tiny, draw=none, fill=white, fill opacity=0.85, text opacity=1, at={(0.98,0.98)}, anchor=north east}, legend columns=2]
          \fill[mGood!9] (axis cs:-0.3,0) rectangle (axis cs:5.3,15);
          \fill[mBad!9] (axis cs:-0.3,-6) rectangle (axis cs:5.3,0);
          \node[text=mGood, font=\scriptsize\bfseries, anchor=north west] at (axis cs:0.9,14.6) {混ぜて得};
          \node[text=mBad, font=\scriptsize\bfseries, anchor=south west] at (axis cs:-0.2,-5.7) {混ぜて損};
          \draw[black, line width=0.7pt] (axis cs:-0.3,0) -- (axis cs:5.3,0);
          \addplot[sA] table[x=x, y=a] {gain_mean.csv}; \addlegendentry{A}
          \addplot[sV] table[x=x, y=v] {gain_mean.csv}; \addlegendentry{V}
          \addplot[sT] table[x=x, y=t] {gain_mean.csv}; \addlegendentry{T}
          \addplot[sAV] table[x=x, y=av] {gain_mean.csv}; \addlegendentry{AV}
          \addplot[sVT] table[x=x, y=vt] {gain_mean.csv}; \addlegendentry{VT}
          \addplot[sAT] table[x=x, y=at] {gain_mean.csv}; \addlegendentry{AT}
          \addplot[sAVT] table[x=x, y=avt] {gain_mean.csv}; \addlegendentry{AVT}
          \addplot[sAH] table[x=x, y=ah] {gain_mean.csv}; \addlegendentry{A(HuBERT)}
        \end{axis}
      \end{tikzpicture}\par
      {\raggedright\tiny 利得の組の平均．同じコーパスの組を含む全組，各組 30 回．合併はソースのみの $C$ を使う\par}
    \end{column}
    \begin{column}{0.45\textwidth}
      \begin{itemize}\footnotesize\setlength{\itemsep}{10pt}
        \item A の $K=1$ で利得 \NUM{map-gain-a-k1} \pt．\NUM{map-pairs-a} 組のうち \NUM{map-gain-pos-a-k1} 組で区間が 0 より上
        \item $K=25$ では 8 種類すべてで負（A は \NUM{map-gain-a-k25} \pt）
        \item 8 種類とも同じ形をとる
      \end{itemize}
    \end{column}
  \end{columns}
  \hosoku{s:p1:map}{種類ごとの値}
\end{frame}

\begin{frame}{得の出る組は $K$ とともに減り，得の大きさは主に 2 つのコーパスの性質で決まる}
  \rl{結果：利得の総当たり実験}
  \vspace{-0.3em}
  \centering
  \begin{tikzpicture}
    \begin{axis}[width=14.2cm, height=4.6cm, ybar stacked, bar width=9pt, ymin=0, ymax=1, xmin=-0.7, xmax=26.7,
        xtick={0,...,5,7,8,...,12,14,15,...,19,21,22,...,26},
        xticklabels={1,5,10,25,50,全,1,5,10,25,50,全,1,5,10,25,50,全,1,5,10,25,50,全},
        xticklabel style={font=\tiny}, ytick={0,0.25,0.5,0.75,1}, ylabel={組の割合},
        legend style={font=\tiny, draw=none, at={(1,1.02)}, anchor=south east, legend columns=3, /tikz/every even column/.append style={column sep=0.8em}},
        clip=false]
      \addplot[fill=cOurs, draw=none] table[x=x, y=pos] {gain_sign.csv}; \addlegendentry{95\% 区間が 0 を上回る}
      \addplot[fill=kAV!35, draw=none] table[x=x, y=mid] {gain_sign.csv}; \addlegendentry{区間が 0 をまたぐ}
      \addplot[fill=cBest, draw=none] table[x=x, y=neg] {gain_sign.csv}; \addlegendentry{区間が 0 を下回る}
      \node[font=\scriptsize, anchor=north] at (axis cs:2.5,-0.13) {A（\NUM{map-pairs-a} 組）};
      \node[font=\scriptsize, anchor=north] at (axis cs:9.5,-0.13) {V（\NUM{map-pairs-v} 組）};
      \node[font=\scriptsize, anchor=north] at (axis cs:16.5,-0.13) {AV（\NUM{map-pairs-v} 組）};
      \node[font=\scriptsize, anchor=north] at (axis cs:23.5,-0.13) {T（\NUM{map-pairs-t} 組）};
    \end{axis}
  \end{tikzpicture}\par
  \vspace{0.5em}\raggedright
  \begin{itemize}\footnotesize\setlength{\itemsep}{4pt}
        \item 区間が 0 を上回る組（青）は $K$ とともに減り，下回る組（橙）が増える
        \item ターゲットとソースの効果を足す模型で，利得のばらつきの \NUM{map-share-a-lo}〜\NUM{map-share-v-hi}\% を説明できる
      \end{itemize}
  \hosoku{s:p1:map}{加法の模型と種類ごとの値}
\end{frame}

\begin{frame}{$K$ 本で推定した判断基準の誤差は $1/(2K)$，向きの誤差は $q/K$ で減る}
  \rl{考え方：2 つのコスト}
  \vspace{-0.2em}
  \begin{columns}[T,onlytextwidth]
    \begin{column}{0.47\textwidth}
      \centering
      \begin{tikzpicture}
        \pgfplotstablegetelem{79}{thr}\of{figures/cost_curves.csv}\edef\cthr{\pgfplotsretval}%
        \pgfplotstablegetelem{79}{a}\of{figures/cost_curves.csv}\edef\cca{\pgfplotsretval}%
        \pgfplotstablegetelem{79}{ah}\of{figures/cost_curves.csv}\edef\ccah{\pgfplotsretval}%
        \pgfplotstablegetelem{79}{av}\of{figures/cost_curves.csv}\edef\ccav{\pgfplotsretval}%
        \pgfplotstablegetelem{79}{v}\of{figures/cost_curves.csv}\edef\ccv{\pgfplotsretval}%
        \begin{loglogaxis}[width=5.6cm, height=5.4cm, xmin=1, xmax=50, ymin=0.005, ymax=40, clip=false,
            xtick={1,5,10,25,50}, xticklabels={1,5,10,25,50}, xlabel={感情あたりのラベル $K$}, ylabel={$K$ 本で推定した誤差},
            ymajorgrids, grid style={mMuted!25}]
          \addplot[black, line width=1.8pt] table[x=K, y=thr] {cost_curves.csv};
          \addplot[sAH, no marks] table[x=K, y=ah] {cost_curves.csv};
          \addplot[sA, no marks] table[x=K, y=a] {cost_curves.csv};
          \addplot[sAV, no marks] table[x=K, y=av] {cost_curves.csv};
          \addplot[sV, no marks] table[x=K, y=v] {cost_curves.csv};
          \node[anchor=west, font=\scriptsize\bfseries] at (axis cs:52,\cthr) {判断基準};
          \node[anchor=west, font=\tiny, text=kAH] at (axis cs:52,\ccah) {\raisebox{0.6em}{向き A(Hu)}};
          \node[anchor=west, font=\tiny, text=kA] at (axis cs:52,\cca) {\raisebox{-0.6em}{向き A}};
          \node[anchor=west, font=\tiny, text=kAV!80!black] at (axis cs:52,\ccav) {向き AV};
          \node[anchor=west, font=\tiny, text=kV] at (axis cs:52,\ccv) {向き V};
          \node[anchor=south west, font=\scriptsize] at (axis cs:1.1,0.007) {判断基準 $1/(2K)$}; \node[anchor=south west, font=\scriptsize, text=cDisc] at (axis cs:7,8) {向き $q/K$};
        \end{loglogaxis}
      \end{tikzpicture}\par
      {\raggedright\tiny 式の値（$q$ はターゲットの中央値）．両軸とも対数．向きの誤差は判断基準の誤差の $2q$ 倍\par}
    \end{column}
    \begin{column}{0.51\textwidth}
      \begin{itemize}\footnotesize\setlength{\itemsep}{10pt}
        \item \crit{判断基準}：$K$ 本で決めた真ん中のずれは $1/(2K)$ で減り，数本で大半が消える
        \item \disc{向き}：識別力は $d(K) = \Delta/\sqrt{1+q/K}$．$q$ は A \NUM{p1-qdir-a}，V \NUM{p1-qdir-v}
        \item 音声で最良に近い向きを得るには，感情あたり 20 本ほど要る
      \end{itemize}
    \end{column}
  \end{columns}
  \hosoku{s:p1:thr}{2 つのコストの導出}
\end{frame}

\begin{frame}{2 つのコストから 4 つの予想が出る}
  \rl{予想}
  \vspace{0.3em}
  \centering\footnotesize
  \renewcommand{\arraystretch}{1.35}
  \begin{tabularx}{\textwidth}{@{}l X X@{}}
    \toprule
    & 予想 & 確かめ方 \\
    \midrule
    P1 & 少ないラベルが取り戻すソースのみの損の大半は，判断基準による損失である & 合併がソースのみを上回った分を，$o$ の変化の分と $d$ の変化の分に分ける \\
    P2 & 上回った分に占める判断基準の割合は，$K$ が増えると減る & 同じ割合を $K = 1, 5, 10$ で比べる \\
    P3 & 合併は，ターゲットが受ける比重に応じて判断基準を直す & Newton 法の 1 歩から直る割合 $\lambda$ を予言し，測った $\lambda$ と比べる \\
    P4 & ソースは，$q$ と，ソースの向きがターゲットにどれだけ合うかで決まる本数のラベルの代わりになる & ソースの価値 $n$ を予言し，学習曲線の上で測った $n$ と比べる \\
    \bottomrule
  \end{tabularx}
\end{frame}

\begin{frame}{ソースのみの分類器は主に向きで損をする．顔だけは 2 つの損がほぼ同じ}
  \rl{結果：損の内訳}
  \vspace{-0.2em}
  \begin{columns}[T,onlytextwidth]
    \begin{column}{0.56\textwidth}
      \centering
      \begin{tikzpicture}
        \begin{axis}[width=7.8cm, height=5.3cm, ybar, bar width=8pt, ymin=0, ymax=13, xmin=-0.6, xmax=7.6,
            xtick={0,...,7}, xticklabels={A,V,T,AV,VT,AT,AVT,{A(Hu)}}, xticklabel style={font=\tiny},
            ylabel={損失（対の正解率の\pt）}, ymajorgrids, grid style={mMuted!25},
            legend style={font=\tiny, draw=none, at={(0.5,1.02)}, anchor=south, legend columns=2, /tikz/every even column/.append style={column sep=1em}}]
          \addplot[fill=cCrit, draw=none] table[x=x, y=thr] {loss_split.csv}; \addlegendentry{判断基準による損失}
          \addplot[fill=white, draw=cDisc, pattern=north east lines, pattern color=cDisc] table[x=x, y=dir] {loss_split.csv}; \addlegendentry{向きによる損失}
        \end{axis}
      \end{tikzpicture}\par
      {\raggedright\tiny ターゲットで最良の線形の規則と比べた損失．6 つの感情の対と，違うコーパスの組で平均（A(HuBERT) は全組）\par}
    \end{column}
    \begin{column}{0.42\textwidth}
      \begin{itemize}\footnotesize\setlength{\itemsep}{10pt}
        \item A：判断基準 \NUM{p1-e13-loss-offset-a-xsame} に対し向き \NUM{p1-e13-loss-dir-a-xsame} \pt
        \item V：\NUM{p1-e13-loss-offset-v-xsame} と \NUM{p1-e13-loss-dir-v-xsame} \pt でほぼ同じ
        \item T：\NUM{p1-e13-loss-offset-t-xsame} と \NUM{p1-e13-loss-dir-t-xsame} \pt．判断基準の損が小さい
      \end{itemize}
    \end{column}
  \end{columns}
  \hosoku{s:p1:loss}{損の分け方と全部の値}
\end{frame}

\begin{frame}{判断基準の損の多くはコーパスの平均のずれで，ラベル無しの中心化で消える}
  \rl{結果：中心化}
  \vspace{-0.2em}
  \begin{columns}[T,onlytextwidth]
    \begin{column}{0.55\textwidth}
      \centering
      \pgfplotstablegetelem{0}{a_center}\of{figures/centering_t.csv}\edef\cca{\pgfplotsretval}%
      \pgfplotstablegetelem{0}{ah_center}\of{figures/centering_t.csv}\edef\ccah{\pgfplotsretval}%
      \pgfplotstablegetelem{0}{v_center}\of{figures/centering_t.csv}\edef\ccv{\pgfplotsretval}%
      \begin{tikzpicture}
        \begin{axis}[width=7.8cm, height=5.2cm, ybar, bar width=6.5pt, ymin=-2, ymax=6.5, xmin=-0.5, xmax=4.5,
            xtick={0,...,4}, xticklabels={\shortstack{中心化\\（ラベル無し）},\shortstack{切片だけ\\学び直す},合併,\shortstack{中心化\\＋切片},\shortstack{中心化\\＋合併}},
            xticklabel style={font=\tiny, align=center}, ylabel={ソースのみからの差（\pt）},
            ymajorgrids, grid style={mMuted!25}, legend style={font=\tiny, draw=none, at={(0.5,1.02)}, anchor=south, legend columns=4},
            clip=false]
          \draw[black, line width=0.6pt] (axis cs:-0.5,0) -- (axis cs:4.5,0);
          \draw[mMuted, densely dashed] (axis cs:2.5,-2) -- (axis cs:2.5,6.5);
          \addplot[fill=kA, draw=none] table[x=x, y=a] {centering_t.csv}; \addlegendentry{A}
          \addplot[fill=kAH, draw=none] table[x=x, y=ah] {centering_t.csv}; \addlegendentry{A(HuBERT)}
          \addplot[fill=kV, draw=none] table[x=x, y=v] {centering_t.csv}; \addlegendentry{V}
          \addlegendimage{line legend, black, line width=1pt}\addlegendentry{中心化だけの高さ}
          % 中心化だけの高さ（後ろの 2 群）．棒の中心は x + (−1, 0, +1) × 8.5pt（棒の幅 6.5pt ＋ 間 2pt）
          % \foreach は使わない（pgfplots は軸の中の描画を後回しにするので，ループの変数が消える）
          \draw[black, line width=1.0pt] (axis cs:3,\cca) ++(-11.75pt,0) -- ++(6.5pt,0);
          \draw[black, line width=1.0pt] (axis cs:3,\ccah) ++(-3.25pt,0) -- ++(6.5pt,0);
          \draw[black, line width=1.0pt] (axis cs:3,\ccv) ++(5.25pt,0) -- ++(6.5pt,0);
          \draw[black, line width=1.0pt] (axis cs:4,\cca) ++(-11.75pt,0) -- ++(6.5pt,0);
          \draw[black, line width=1.0pt] (axis cs:4,\ccah) ++(-3.25pt,0) -- ++(6.5pt,0);
          \draw[black, line width=1.0pt] (axis cs:4,\ccv) ++(5.25pt,0) -- ++(6.5pt,0);
          \node[font=\tiny, text=mMuted] at (axis cs:1.5,6.1) {$K=5$};
        \end{axis}
      \end{tikzpicture}\par
      {\raggedright\tiny 組の平均．中心化：コーパスごとにラベル無しの発話の平均を引く．合併は総当たり実験の値，中心化の版は学び直し（$C$ は交差検証，反復 3 回）\par}
    \end{column}
    \begin{column}{0.43\textwidth}
      \begin{itemize}\footnotesize\setlength{\itemsep}{10pt}
        \item 中心化だけで，ソースのみの UAR が A \NUM{p1-center-srconly-a}，V \NUM{p1-center-srconly-v} \pt 上がる
        \item $K=5$ 本で切片だけを学び直すと A で \NUM{p1-d1-gain-a-k5} \pt 上がり，合併（\NUM{p1-merged-gain-a-k5} \pt）に近い
        \item 中心化の後に足す分は A で \NUM{p1-center-d1-gain-a-k5}（切片）と \NUM{p1-center-merged-gain-a-k5} \pt（合併）
      \end{itemize}
    \end{column}
  \end{columns}
  \hosoku{s:p1:center}{中心化と切片の学び直しの詳細}
\end{frame}

\begin{frame}{単独の A・V の $K\le5$ では，上乗せの \NUM{p1-e29-critshare-kle5-single-lo}〜\NUM{p1-e29-critshare-kle5-single-hi} が判断基準の分}
  \rl{結果：P1・P2}
  \vspace{-0.2em}
  \begin{columns}[T,onlytextwidth]
    \begin{column}{0.47\textwidth}
      \centering
      \begin{tikzpicture}
        \pgfplotstablegetelem{2}{a}\of{figures/thrshare.csv}\edef\tsa{\pgfplotsretval}%
        \pgfplotstablegetelem{2}{ah}\of{figures/thrshare.csv}\edef\tsah{\pgfplotsretval}%
        \pgfplotstablegetelem{2}{v}\of{figures/thrshare.csv}\edef\tsv{\pgfplotsretval}%
        \pgfplotstablegetelem{2}{t}\of{figures/thrshare.csv}\edef\tst{\pgfplotsretval}%
        \begin{axis}[width=5.8cm, height=5.4cm, xmin=-0.15, xmax=2.15, ymin=0, ymax=1, clip=false,
            xtick={0,1,2}, xticklabels={1,5,10}, xlabel={感情あたりのラベル $K$},
            ylabel={上乗せに占める判断基準の割合}, ymajorgrids, grid style={mMuted!25}]
          \draw[black, line width=0.7pt] (axis cs:-0.15,0.5) -- (axis cs:2.15,0.5);
          \node[anchor=south west, font=\tiny] at (axis cs:-0.12,0.5) {半分};
          \addplot[mMuted!70, dashed, line width=0.8pt, mark=none] table[x=x, y=av] {thrshare.csv};
          \addplot[mMuted!70, dashed, line width=0.8pt, mark=none] table[x=x, y=vt] {thrshare.csv};
          \addplot[mMuted!70, dashed, line width=0.8pt, mark=none] table[x=x, y=avt] {thrshare.csv};
          \addplot[mMuted!70, dashed, line width=0.8pt, mark=none] table[x=x, y=at] {thrshare.csv};
          \addplot[mMuted!70, dashed, line width=0.8pt, mark=none] table[x=x, y=t] {thrshare.csv};
          \addplot[sA, line width=1.6pt] table[x=x, y=a] {thrshare.csv};
          \addplot[sAH, line width=1.6pt] table[x=x, y=ah] {thrshare.csv};
          \addplot[sV, line width=1.6pt] table[x=x, y=v] {thrshare.csv};
          \node[anchor=west, font=\scriptsize, text=kV] at (axis cs:2.17,\tsv) {\raisebox{0.4em}{V}};
          \node[anchor=west, font=\scriptsize, text=kA] at (axis cs:2.17,\tsa) {A};
          \node[anchor=west, font=\scriptsize, text=kAH] at (axis cs:2.17,\tsah) {\raisebox{0.1em}{A(Hu)}};
          \node[anchor=south west, font=\tiny, text=mMuted] at (axis cs:-0.1,0.03) {灰の破線：融合と言葉を含む 5 種類};
        \end{axis}
      \end{tikzpicture}\par
      {\raggedright\tiny 上乗せ：合併がソースのみを上回った分．組の中央値（IEMOCAP の 2 群の組を除く，反復 10 回）\par}
    \end{column}
    \begin{column}{0.51\textwidth}
      \begin{itemize}\footnotesize\setlength{\itemsep}{10pt}
        \item 損に占める割合（\NUM{p1-e29-lossshare-lo}〜\NUM{p1-e29-lossshare-hi}）よりずっと大きい：ラベルは安い部分を先に直す（P1）
        \item $K$ とともに減る（P2）
        \item 融合と言葉の $K=5$ では半分前後以下（\NUM{p1-e29-critshare-k5-fused-text-lo}〜\NUM{p1-e29-critshare-k5-fused-text-hi}）
      \end{itemize}
    \end{column}
  \end{columns}
  \hosoku{s:p1:share}{割合の求め方と全部の値}
\end{frame}

\begin{frame}{判断基準の割合は，判断基準の損が大きい音声と顔で大きい（論文の表 2）}
  \rl{結果：まとめの表}
  \vspace{0.2em}
  \centering\scriptsize
  \setlength{\tabcolsep}{7pt}
  \renewcommand{\arraystretch}{1.1}
  \begin{tabular}{@{}lrrrrrr@{}}
    \toprule
     & \multicolumn{2}{c}{ソースのみの損失（\pt）} & & \multicolumn{3}{c}{上乗せに占める判断基準の割合} \\
    \cmidrule(lr){2-3}\cmidrule(l){5-7}
    種類 & \crit{判断基準} & \disc{向き} & $q$ & $K{=}1$ & $K{=}5$ & $K{=}10$ \\
    \midrule
    A & \NUM{p1-e13-loss-offset-a-xsame} & \NUM{p1-e13-loss-dir-a-xsame} & \NUM{p1-qdir-a} & \NUM{p1-e29-critshare-k1-a} & \NUM{p1-e29-critshare-k5-a} & \NUM{p1-e29-critshare-k10-a} \\
    A (HuBERT) & \NUM{p1-e13-loss-offset-ah} & \NUM{p1-e13-loss-dir-ah} & \NUM{p1-qdir-ah} & \NUM{p1-e29-critshare-k1-ah} & \NUM{p1-e29-critshare-k5-ah} & \NUM{p1-e29-critshare-k10-ah} \\
    V & \NUM{p1-e13-loss-offset-v-xsame} & \NUM{p1-e13-loss-dir-v-xsame} & \NUM{p1-qdir-v} & \NUM{p1-e29-critshare-k1-v} & \NUM{p1-e29-critshare-k5-v} & \NUM{p1-e29-critshare-k10-v} \\
    AV & \NUM{p1-e13-loss-offset-av-xsame} & \NUM{p1-e13-loss-dir-av-xsame} & \NUM{p1-qdir-av} & \NUM{p1-e29-critshare-k1-av} & \NUM{p1-e29-critshare-k5-av} & \NUM{p1-e29-critshare-k10-av} \\
    T & \NUM{p1-e13-loss-offset-t-xsame} & \NUM{p1-e13-loss-dir-t-xsame} & -- & \NUM{p1-e29-critshare-k1-t} & \NUM{p1-e29-critshare-k5-t} & \NUM{p1-e29-critshare-k10-t} \\
    VT & \NUM{p1-e13-loss-offset-vt-xsame} & \NUM{p1-e13-loss-dir-vt-xsame} & -- & \NUM{p1-e29-critshare-k1-vt} & \NUM{p1-e29-critshare-k5-vt} & \NUM{p1-e29-critshare-k10-vt} \\
    AT & \NUM{p1-e13-loss-offset-at-xsame} & \NUM{p1-e13-loss-dir-at-xsame} & -- & \NUM{p1-e29-critshare-k1-at} & \NUM{p1-e29-critshare-k5-at} & \NUM{p1-e29-critshare-k10-at} \\
    AVT & \NUM{p1-e13-loss-offset-avt-xsame} & \NUM{p1-e13-loss-dir-avt-xsame} & -- & \NUM{p1-e29-critshare-k1-avt} & \NUM{p1-e29-critshare-k5-avt} & \NUM{p1-e29-critshare-k10-avt} \\
    \bottomrule
  \end{tabular}\par
  \vspace{0.4em}
  {\raggedright\tiny 違うコーパスの組で平均（A(HuBERT) の損失は全組）．割合は組の中央値．$q$ は言葉を含む種類では求めていない\par}
  \vspace{0.3em}
  \takeaway{言葉は判断基準の損が小さいので，ラベルが直す判断基準も少ない：上乗せの割合は損の割合に従う}
\end{frame}

\begin{frame}{音声の向きは，ソースが残した所からほとんど動かない}
  \rl{結果：向き}
  \vspace{-0.2em}
  \begin{columns}[T,onlytextwidth]
    \begin{column}{0.5\textwidth}
      \centering
      \begin{tikzpicture}
        \begin{axis}[width=6.9cm, height=5.0cm, ybar, bar width=10pt, ymin=0, ymax=100, xmin=-0.6, xmax=3.6,
            xtick={0,...,3}, xticklabels={A,A(HuBERT),AV,V}, ylabel={取り戻した向きの損失（\%）},
            ymajorgrids, grid style={mMuted!25}, legend style={font=\tiny, draw=none, at={(0.02,0.98)}, anchor=north west}]
          \addplot[fill=white, draw=cDisc, pattern=north east lines, pattern color=cDisc] table[x=x, y=bound] {dirrecov.csv}; \addlegendentry{式の上限（最適な組み合わせ）}
          \addplot[fill=cDisc, draw=none] table[x=x, y=ht] {dirrecov.csv}; \addlegendentry{仮説転移の実測}
        \end{axis}
      \end{tikzpicture}\par
      {\raggedright\tiny $K=10$．組の中央値（違うコーパスの組）\par}
    \end{column}
    \begin{column}{0.48\textwidth}
      \begin{itemize}\footnotesize\setlength{\itemsep}{10pt}
        \item $K=10$ 本と最適に組み合わせても，音声の向きの損の \NUM{p1-dirloss-recov-bound-k10-audio-lo}〜\NUM{p1-dirloss-recov-bound-k10-audio-hi}\% しか戻らない（式の上限）
        \item 仮説転移の実測は A で \NUM{p1-dirloss-recov-ht-k10-a}\%，A(HuBERT) で \NUM{p1-dirloss-recov-ht-k10-ah}\%
      \end{itemize}
    \end{column}
  \end{columns}
  \hosoku{s:p1:dirrec}{上限の式と全部の値}
\end{frame}

\begin{frame}{合併が判断基準を直す割合 $\lambda$ は，1 つの量 $\rho$ で予言できる}
  \rl{結果：P3}
  \vspace{-0.2em}
  \begin{columns}[T,onlytextwidth]
    \begin{column}{0.52\textwidth}
      \begin{itemize}\scriptsize\setlength{\itemsep}{7pt}
        \item Newton 法の 1 歩から $\lambda = \rho/(1+\rho)$ を導いた．$\rho$ は 2 つの経路の和
        \item 切片の経路：ターゲットが学習の比重に占める割合
        \item 重みの経路：ソースの散らばらない向きにターゲットが離れていると，その向きの重みがターゲットだけの切片になる
        \item 予言と実測の順位相関は \NUM{p1-e15-lam-rho-lo}〜\NUM{p1-e15-lam-rho-hi}（平均は大きめに予言する）
      \end{itemize}
    \end{column}
    \begin{column}{0.46\textwidth}
      \centering
      \begin{tikzpicture}
        \begin{axis}[width=5.3cm, height=5.0cm, xmin=-0.22, xmax=1.06, ymin=-0.22, ymax=1.06,
            xlabel={予言した $\lambda = \rho/(1+\rho)$}, ylabel={合併で測った $\lambda$},
            xtick={0,0.25,0.5,0.75,1}, ytick={0,0.25,0.5,0.75,1}, ymajorgrids, xmajorgrids, grid style={mMuted!20},
            legend style={font=\tiny, draw=none, fill=none, at={(0.5,-0.3)}, anchor=north, legend columns=2, /tikz/every even column/.append style={column sep=0.6em}}]
          \draw[mMuted, densely dashed] (axis cs:-0.22,-0.22) -- (axis cs:1.06,1.06);
          \addplot[only marks, mark=*, mark size=1.0pt, kA, opacity=0.7] table[x=pred, y=obs] {lambda_a.csv}; \addlegendentry{A：\NUM{p1-e15-lam-rho-a}}
          \addplot[only marks, mark=star, mark size=1.4pt, kAH, opacity=0.7] table[x=pred, y=obs] {lambda_ah.csv}; \addlegendentry{A(HuBERT)：\NUM{p1-e15-lam-rho-ah}}
          \addplot[only marks, mark=square*, mark size=1.0pt, kV, opacity=0.7] table[x=pred, y=obs] {lambda_v.csv}; \addlegendentry{V：\NUM{p1-e15-lam-rho-v}}
          \addplot[only marks, mark=diamond*, mark size=1.3pt, kAV, opacity=0.7] table[x=pred, y=obs] {lambda_av.csv}; \addlegendentry{AV：\NUM{p1-e15-lam-rho-av}}
        \end{axis}
      \end{tikzpicture}\par
      {\raggedright\tiny 点は組 $\times$ $K \in \{1,5,10\}$（反復 10 回）．凡例の値は順位相関．破線は $y=x$\par}
    \end{column}
  \end{columns}
  \hosoku{s:p1:lam}{$\lambda$ の導出}
\end{frame}

\begin{frame}{ソースから出発すると，ターゲットだけで $K+n$ 本学んだのと同じになる}
  \rl{結果：P4　ソースの価値}
  \vspace{-0.2em}
  \begin{columns}[T,onlytextwidth]
    \begin{column}{0.46\textwidth}
      \centering
      \begin{tikzpicture}
        \begin{semilogxaxis}[width=6.4cm, height=5.0cm, xmin=1, xmax=100, ymin=55, ymax=92,
            xtick={1,5,15,50,100}, xticklabels={1,$K$,$K+n$,50,100}, ytick=\empty,
            xlabel={ターゲットだけで学ぶラベルの本数}, ylabel={正解率}, clip=false]
          \addplot[cOurs, line width=1.3pt] table[x=K, y=acc] {value_curve.csv};
          \node[font=\tiny, text=cOurs, anchor=west] at (axis cs:22,84) {ターゲットのみの学習曲線};
          \addplot[only marks, mark=*, mark size=2pt, cOurs] table[x=K, y=acc] {value_points.csv};
          \pgfplotstablegetelem{1}{acc}\of{figures/value_points.csv}\edef\accn{\pgfplotsretval}
          \addplot[only marks, mark=*, mark size=2.4pt, cExt] coordinates {(5,\accn)};
          \node[font=\tiny, text=cExt, anchor=south east] at (axis cs:5,\accn) {仮説転移（$K$ 本）};
          \draw[->, cExt, line width=0.9pt] (axis cs:5.6,\accn) -- (axis cs:13.8,\accn);
          \node[font=\tiny, text=cExt, anchor=north] at (axis cs:8.8,\accn) {$n$};
        \end{semilogxaxis}
      \end{tikzpicture}\par
      {\raggedright\tiny 模式図（A の $q$ と説明用の $\Delta$，$n=10$ で描いた）\par}
    \end{column}
    \begin{column}{0.52\textwidth}
      \begin{itemize}\footnotesize\setlength{\itemsep}{10pt}
        \item $e$：ソースの向きの効率（$d = e\Delta$）．$n = q e^2/(1-e^2)$ で，$n$ は $K$ によらない
        \item 1 ソースの音声で \NUM{p1-e32-n1src-audio-lo}〜\NUM{p1-e32-n1src-audio-hi} 本．全ソースの A で $K=5$ に \NUM{p1-e32-nall-a-k5} 本
        \item 音声の組で予言と実測の順位相関は \NUM{p1-e32-rho-audio-lo}〜\NUM{p1-e32-rho-audio-hi}
      \end{itemize}
    \end{column}
  \end{columns}
  \hosoku{s:p1:val}{$n$ の導出と散布図}
\end{frame}

\begin{frame}{第 I 部のまとめ：ラベルは安い判断基準を先に直し，向きは直りにくい}
  \rl{まとめ}
  \vspace{0.2em}
  \centering\scriptsize
  \renewcommand{\arraystretch}{1.3}
  \begin{tabularx}{\textwidth}{@{}l X X@{}}
    \toprule
    & 予想 & 測った結果 \\
    \midrule
    P1 & 少ないラベルが取り戻す損の大半は判断基準 & 単独の A・V の $K\le5$ で上乗せの \NUM{p1-e29-critshare-kle5-single-lo}〜\NUM{p1-e29-critshare-kle5-single-hi} が判断基準の分．融合と言葉の $K=5$ では半分前後以下 \\
    P2 & 判断基準の割合は $K$ とともに減る & 8 種類とも $K = 1 \to 10$ で減る \\
    P3 & 合併はターゲットの比重に応じて判断基準を直す & 予言と測った $\lambda$ の順位相関 \NUM{p1-e15-lam-rho-lo}〜\NUM{p1-e15-lam-rho-hi}．平均は大きめに予言する \\
    P4 & ソースは $n = q e^2/(1-e^2)$ 本ぶんのラベルになる & 音声の 1 ソースで \NUM{p1-e32-n1src-audio-lo}〜\NUM{p1-e32-n1src-audio-hi} 本．順位相関 \NUM{p1-e32-rho-audio-lo}〜\NUM{p1-e32-rho-audio-hi} \\
    \bottomrule
  \end{tabularx}\par
  \vspace{0.5em}
  \takeaway{判断基準は中心化と切片の学び直しで足りる．残る主な損は向きで，音声では $q$ 本ほどのラベルが要り，ソース 1 つは約 10 本ぶんしか補わない}
\end{frame}
